\section{Two column constraints along a path of principal minors}\label{sec:paths}

Proposition~\ref{prop:static-column} bounds the admissible columns when
\(W_2\) is small. We now show that a terminal matrix with large \(W_2\)
must encounter two rare column constraints along a short path of principal
minors. Each constraint costs almost \(b^N\), giving the exceptional
probability required in Proposition~\ref{prop:column-bound}.
The principal-minor method and its layer recursion come from
\cite[Sections 7 and 11]{SymmetricRate}. We first state the deterministic
inputs, then give the probability estimates for biased columns.

\subsection{Principal extensions and retained layers}

A principal extension of width \(h\in\{1,2\}\) has the form
\begin{equation}\label{eq:path-extension}
 B=\begin{pmatrix}C&V\\V^{\mathsf T}&D\end{pmatrix},
 \qquad V=(v_1\ \cdots\ v_h),
\end{equation}
where \(C\) is nonsingular and \(D\) is symmetric. At width two we impose
the \emph{gate}
\begin{equation}\label{eq:path-gate}
 v_1,v_2\in\mathcal S_{L_0}(C).
\end{equation}
There is no gate at width one. The following elementary fact ensures that
every nonsingular bit matrix can be reached by a path with these gates.

\begin{lemma}[Gated deletion, {\cite[Lemma 7.1]{SymmetricRate}}]
\label{lem:path-deletion}
Let \(B\) be a nonsingular symmetric bit matrix of positive size.
The empty principal minor has determinant one, and its admissible
column set consists of the unique empty column.
Either some one-coordinate principal deletion is nonsingular, or a
two-coordinate principal deletion gives a nonsingular matrix \(C\) for
which both deleted cross columns satisfy
\[
 v_i^{\mathsf T}C^{-1}v_i\in\bits,
 \qquad \lVert C^{-1}v_i\rVert_\infty\le1
 \quad(i=1,2).
\]
In particular, whenever the sizes belong to \(\mathcal I_N\), the
two-coordinate move satisfies Equation~\eqref{eq:path-gate}.
\end{lemma}

\begin{proof}
The cofactor identity gives the first alternative if some diagonal entry
of \(B^{-1}\) is nonzero. Otherwise choose distinct \(i,j\) maximizing
the absolute value of an entry of \(B^{-1}\), and write
\(\kappa=(B^{-1})_{ij}\ne0\). The corresponding principal block is
\(\left(\begin{smallmatrix}0&\kappa\\\kappa&0\end{smallmatrix}\right)\).
Jacobi's complementary minor identity makes the complementary matrix
\(C\) nonsingular. In the block notation of
Equation~\eqref{eq:path-extension}, its Schur complement
\(T=D-V^{\mathsf T}C^{-1}V\) satisfies
\[
 T^{-1}=\begin{pmatrix}0&\kappa\\\kappa&0\end{pmatrix},
 \qquad (B^{-1})_{I,\{i,j\}}=-C^{-1}VT^{-1},
\]
where \(I\) is the set of remaining indices. The diagonal of \(T\)
vanishes, giving the two quadratic identities. The two columns of the
inverse block on the right are \(-\kappa C^{-1}v_2\) and
\(-\kappa C^{-1}v_1\). Their entries have absolute value at most
\(|\kappa|\), proving the norm bound. Finally
\(\lVert C^{-1}v_i\rVert_2\le\sqrt{\operatorname{size}(C)}
\le\sqrt N\le L_0\sqrt N\).
\end{proof}

For an extension with nonsingular start \(C\), define the successive
quotients
\begin{equation}\label{eq:path-quotients}
 G_0=G_C,\qquad
 G_i=\Z^m/(C\Z^m+\Z v_1+\cdots+\Z v_i)
       =G_{i-1}/\langle\overline v_i\rangle
       \quad(1\le i\le h),
\end{equation}
where \(m=\operatorname{size}(C)\). These groups are defined even if
the endpoint \(B\) is singular. In a one-element quotient, each layer
height drops by zero or one. A positive layer whose height does not
drop is called \emph{retained}. A \emph{plateau} is a maximal interval
of consecutive layers of equal positive height at one prime. On each
plateau, the retained layers form a prefix, possibly empty. These
facts are Lemma~7.3 of \cite{SymmetricRate}; for example the prefix
assertion follows because quotient heights are nonincreasing and can
only equal the original height or one less.

Here are the remaining group-theoretic inputs in the precise form
needed below.

\begin{lemma}[Retention lattices and tail recursion]
\label{lem:path-algebra}
Let \(C\) be a nonsingular symmetric bit matrix of size \(m\le N\).
\begin{enumerate}[label=\textnormal{(\roman*)}]
\item Let \(\pi:\Z^m\twoheadrightarrow G\) be a quotient map with
\(C\Z^m\subseteq\ker\pi\). For any set \(\mathcal D\) of positive
layers of \(G\), there is a lattice \(K_{\mathcal D}\), fixed by
\(\pi,G,\mathcal D\), such that
\begin{equation}\label{eq:path-retention-lattice}
 C\Z^m\subseteq K_{\mathcal D}\subseteq\Z^m,
 \qquad [\Z^m:K_{\mathcal D}]=N^{Nq(\mathcal D)},
 \quad q(\mathcal D)=\sum_{(l,j)\in\mathcal D}k_G(l,j)\nu_l.
\end{equation}
If every layer in \(\mathcal D\) is retained on quotienting by
\(\pi(v)\), then \(v\in K_{\mathcal D}\).
\item Suppose the extension in Equation~\eqref{eq:path-extension}
also has nonsingular endpoint \(B\). Then
\(k_B(l,j)\le k_C(l,j)+h\) for every layer. Writing
\(\mathcal R_i\) for the retained layers in \(G_{i-1}\to G_i\),
for every integer \(s\ge h\) we have
\begin{equation}\label{eq:path-tail-recursion}
 W_s(B)\le W_{s+h}(C)
   +2\sum_{i=1}^h
       \sum_{\substack{(l,j)\in\mathcal R_i\\
                        k_{G_{i-1}}(l,j)\ge s+1-h}}\nu_l.
\end{equation}
\end{enumerate}
\end{lemma}

\begin{proof}[Source and interpretation]
Part~\textnormal{(i)} is the Retention lattice lemma
\cite[Lemma 7.4]{SymmetricRate}, including arbitrary, nonconsecutive
sets of layers. Part~\textnormal{(ii)} is the Tail recursion lemma
\cite[Lemma 11.1]{SymmetricRate} together with common-quotient
interlacing \cite[Equation (7.11)]{SymmetricRate}. For the latter,
projection onto the old coordinates identifies \(G_h\) with the
quotient of \(G_B\) by the \(h\) new coordinate classes. Since a
one-element quotient lowers each height by at most one,
\(k_B\le k_{G_h}+h\le k_C+h\).
The stronger recursion uses symmetry through the Symmetric kernel
inequality \cite[Lemma 7.5]{SymmetricRate}; it needs no assumption on
the entry distribution. All of these statements therefore apply to
the present biased model.
\end{proof}

\subsection{The cost of retaining layers}

Fix \(L_0\ge1\), a sufficiently small fixed \(e_0>0\), and the
corresponding \(\zeta>0\) in Proposition~\ref{prop:static-column}.
Recall \(R=\lceil\ell^2\rceil\) and \(\sigma=R^{-6}\).
Call an extension \emph{costly} if, in at least one of its quotient
operations,
\begin{equation}\label{eq:path-costly-definition}
 \sum_{(l,j)\in\mathcal R_i}\nu_l\ge\sigma.
\end{equation}
We use this name only for the condition in
Equation~\eqref{eq:path-costly-definition}.

In all column estimates, the information already exposed fixes the
start \(C\), and possibly a preceding column of the same width-two
extension. The next column has independent bit entries with success
probabilities \(p\) or \(b\). The quotient map is thus fixed before
the next column is tested.

\begin{lemma}[Almost full cost of retention]\label{lem:path-retention-cost}
Typically for nonsingular \(C\) with size in \(\mathcal I_N\),
conditional on \(C\), the probability that an independent principal
extension of either width is costly is at most
\begin{equation}\label{eq:path-retention-cost}
 p_N:=2^{-N(a-2e_0)}.
\end{equation}
For one specified column, the probability that its retained weight
is at least \(\sigma\) obeys the same bound, also after exposing the
preceding column of a width-two extension. Also, typically, when
\(W_2(C)\le\zeta\), the width-two gate has conditional probability
at most \(p_N^2\).
\end{lemma}

\begin{proof}
Impose Theorem~\ref{thm:small-intersection} on \(C\) with
\(\tau=\ell^{-20}\) and some fixed \(0<\alpha<e_0\).
Then \(\tau=o(\sigma)\). We show that the retention condition can be
covered by \(2^{o(N)}\) selected layer sets, each giving a lattice of
index at least \(N^{\sigma N}\).

Let \(G\) be the group before one column is tested. It is a quotient
of \(G_C\), so its total height times weight is at most \(1/2\).
At a prime \(l\), let \(d_l\) be its number of plateaus. Distinct
positive plateau heights give
\(d_l^2\le2\sum_j k_G(l,j)\), and the group-order bound gives
\(\sum_{l,j}k_G(l,j)=O(N\log N)\). Consequently the number of
plateaus at primes \(l\le\sqrt N\) is
\(O(N^{3/4}\sqrt{\log N})\), by Cauchy--Schwarz. Each prefix
length has \(O(N\log N)\) possible values. Thus the complete
retention pattern at these primes can be specified with
\begin{equation}\label{eq:path-small-prime-entropy}
 O\bigl(N^{3/4}(\log N)^{3/2}\bigr)=o(N)
\end{equation}
bits. Include all those retained layers in the selected set.

If their total weight is below \(\sigma\), add retained layers at
primes \(l>\sqrt N\) until the total first reaches \(\sigma\).
There are at most \(N\) candidate layers in this prime range, since
each has \(\nu_l>1/(2N)\). At most
\(\lceil2\sigma N\rceil+1\) additions suffice. The binomial bound
therefore contributes \(2^{o(N)}\) further possibilities. This
constructs the claimed cover whenever the retained weight is at
least \(\sigma\).

For every selected set \(\mathcal D\), positive heights and
Equation~\eqref{eq:path-retention-lattice} give
\(q(\mathcal D)\ge\sigma\). The simultaneous cube estimate applies
to \(K_{\mathcal D}\), since it contains \(C\Z^m\), even when
\(G\) depends on an earlier column. Each cube point has probability
at most \(b^m\), where \(m=N-o(N)\). The one-column probability is
therefore at most
\[
 2^{o(N)}2^{\alpha N}b^m
       =2^{-N(a-\alpha-o(1))}.
\]
The slack \(\alpha<e_0\) absorbs a union over at most two columns
and yields Equation~\eqref{eq:path-retention-cost}.

Finally impose Proposition~\ref{prop:static-column} on the start.
When \(W_2(C)\le\zeta\), each column belongs to
\(\mathcal S_{L_0}(C)\) with probability at most \(2^{e_0N}b^m\).
The two cross columns are independent conditional on \(C\), so the
gate probability is at most
\((2^{e_0N}b^m)^2\le p_N^2\).
\end{proof}

We also need a weaker estimate that applies without typicality. Its
role is to prevent a path from carrying substantial weight in very
high layers. The elementary counting argument below is the biased
version of \cite[Lemma 11.2]{SymmetricRate}.

\begin{lemma}[Dissociation and a lattice constraint]\label{lem:path-dissociation}
Let \(C\) be a nonsingular bit matrix of size \(m\le N\), and let
\(C\Z^m\subseteq K\subseteq\Z^m\) have index \(N^{qN}\).
If \(v\) has independent bit coordinates, each with maximum atom
at most \(b\), then
\begin{equation}\label{eq:path-fixed-mask-cost}
 \Prob(v\in K)\le b^{2qN/3}.
\end{equation}
\end{lemma}

\begin{proof}
Write \(g_1,\ldots,g_m\) for the coordinate images in \(\Z^m/K\).
Choose a maximal set of \(j\) of these images whose subset sums
are all distinct. Such a set is called \emph{dissociated}. Once
the other bits are fixed, at most one assignment of these \(j\)
bits gives zero in the quotient, and that assignment has probability
at most \(b^j\).

Maximality says that every \(g_i\) is a combination of the selected
images with coefficients in \(\{-1,0,1\}\): otherwise adding it
would preserve distinct subset sums. Put these coefficients in a
\(j\)-by-\(m\) integer matrix \(M\), using identity columns at the
selected indices. Then \(\rank M=\rank(MC)=j\), and evaluation
of the selected images is a surjection
\(\theta:\Z^j\twoheadrightarrow\Z^m/K\). Every column of \(MC\)
lies in \(\ker\theta\). Choose \(j\) independent such columns;
each has norm at most \(N\sqrt j\). The index of their integer
span bounds the quotient order, so Hadamard's inequality gives
\[
 N^{qN}\le(N\sqrt j)^j\le N^{3j/2}.
\]
Thus \(j\ge2qN/3\), proving the assertion. If \(j=0\), the
quotient is trivial and the same conclusion holds with empty products.
\end{proof}

For one column with pre-test group \(G\), write
\begin{equation}\label{eq:path-heavy-weight}
 Q_{\rm hi}=\sum_{\substack{(l,j)\text{ retained}\\k_G(l,j)\ge R-1}}
                     k_G(l,j)\nu_l.
\end{equation}

\begin{lemma}[Cost of heavy retained layers]\label{lem:path-heavy-cost}
Fix \(0<\eta<1\), and put \(q_{\min}=\eta/(100\log N)\).
There is a constant \(c_p>0\) such that, uniformly for
\(q_{\min}\le q\le1/2\),
\begin{equation}\label{eq:path-heavy-cost}
 \Prob\bigl(Q_{\rm hi}\in[q,2q)\mid\text{previous exposure}\bigr)
       \le2^{-c_pqN}.
\end{equation}
This holds whenever the previous exposure fixes a nonsingular start
\(C\) of size in \(\mathcal I_N\), and a quotient map to \(G\)
whose kernel contains \(C\Z^m\), while leaving the tested bits
independent with success probabilities \(p\) or \(b\).
\end{lemma}

\begin{proof}
For a fixed set \(\mathcal D\) of retained layers, the retention
lattice and Lemma~\ref{lem:path-dissociation} give cost
\(2^{-(2a/3)q(\mathcal D)N}\). Enumerate the sets consisting of
all retained layers of height at least \(R-1\), with their total
height times weight in \([q,2q)\).

For primes at most \(\sqrt N\), prefix lengths on the eligible
plateaus specify the set. Equation~\eqref{eq:path-small-prime-entropy}
costs
\[
 O\bigl(N^{3/4}(\log N)^{3/2}\bigr)
       =o(N/\log N)=o_\eta(qN)
\]
bits uniformly over the stated range. At larger primes each eligible
layer contributes at least \((R-1)/(2N)\). There are therefore
\(O(N/R)\) eligible layers in total, and a set of weight below
\(2q\) contains at most \(O(qN/R)\) of them. The binomial estimate,
or the all-subsets bound when \(q\) is bounded below, gives logarithmic
count
\[
 O\!\left(\frac{qN}{R}\bigl(1+\log(1/q)\bigr)\right)
       =o_\eta(qN).
\]
Indeed \(1+\log(1/q)=O_\eta(\ell)\), \(R\asymp\ell^2\),
and \(qN/R\to\infty\) uniformly, so rounding has no effect.
Every counted set has \(q(\mathcal D)\ge q\). Its probability
cost absorbs both entropy bounds; for instance any fixed
\(0<c_p<2a/3\) works for all sufficiently large \(N\).
\end{proof}

\subsection{Fixed paths and removal of high layers}

We now turn the column estimates into a bound on terminal matrices.
Let
\[
 H_0=2\lceil N^{3/5}\rceil,\qquad m_*=N-1.
\]
Consider all fixed index paths
\begin{equation}\label{eq:path-index-path}
 I_0\subset I_1\subset\cdots\subset I_{H_0}=[m_*],
 \qquad h_t=|I_t\setminus I_{t-1}|\in\{1,2\}.
\end{equation}
Fix an ordering of each pair of new indices. At a backward step there
are at most \(2N^2\) choices, so the number of paths is at most
\begin{equation}\label{eq:path-count}
 (2N^2)^{H_0}=2^{O(N^{3/5}\log N)}=2^{o(N)}.
\end{equation}
Every node size belongs to \(\mathcal I_N\), since at most
\(2H_0=o(N^{7/10})\) indices are deleted.

Let \(\mathcal T_m\) be the set of nonsingular size-\(m\) matrices
satisfying Lemma~\ref{lem:prime-corank},
Proposition~\ref{prop:static-column}, and the simultaneous cube
property used in Lemma~\ref{lem:path-retention-cost}.
Their exceptional probabilities are uniformly \(2^{-\omega(N)}\).
Every fixed principal submatrix remains in the bit model: its
upper-triangular entries are independent with the prescribed biases,
and its mean is a principal submatrix of the original mean, hence
has rank at most two.
A union bound using Equation~\eqref{eq:path-count} shows that the
probability of any nonsingular node on any path failing these
requirements is still \(2^{-\omega(N)}\).

For one fixed path, write \(B_t=B[I_t]\), and let
\(\mathcal A\) be the event that all its nodes belong to their
respective \(\mathcal T_m\) and every width-two gate holds.
Lemma~\ref{lem:path-deletion}, applied backwards \(H_0\) times,
shows that every nonsingular terminal matrix outside the exceptional
event belongs to \(\mathcal A\) for at least one fixed path.

For the probability estimates, expose a fixed path \emph{forwards}:
first \(B_0\), then the cross columns and new symmetric block of
each extension. Conditional on everything already exposed, a cross
column has independent entries of the original prescribed biases.
The second cross column of a pair remains independent after the first
has been revealed. Future typicality, nonsingularity and gate
requirements enter only through the event \(\mathcal A\).

Call the last \(5R\) steps, together with their \(5R+1\) nodes,
the \emph{terminal window}. We first remove paths with a large
\(W_R\) anywhere in that window.

\begin{lemma}[Small high-layer tails on the window]\label{lem:path-high-tail}
For every fixed \(0<\eta<1\), each fixed path satisfies
\begin{equation}\label{eq:path-high-tail}
 \Prob\!\left(\mathcal A,
       \max_{H_0-5R\le t\le H_0} W_R(B_t)>\eta\right)
       \le2^{-\Omega_{p,\eta}(NR)}.
\end{equation}
\end{lemma}

\begin{proof}
At each tested column form \(Q_{\rm hi}\) from
Equation~\eqref{eq:path-heavy-weight}, and put
\[
 \widetilde Q=
 \begin{cases}
 Q_{\rm hi},&Q_{\rm hi}\ge q_{\min},\\
 0,&Q_{\rm hi}<q_{\min},
 \end{cases}
 \qquad q_{\min}=\frac{\eta}{100\log N}.
\]
Set \(\widetilde Q=0\) for every cross column of a step with singular
start. This defines the quantities on the entire exposure space.
At nonsingular starts, \(Q_{\rm hi}\le1/2\). For dyadic values
\(q_j=2^jq_{\min}\le1/2\), Lemma~\ref{lem:path-heavy-cost} gives,
with any fixed \(0<c'_p<c_p/2\),
\begin{align}
 \E\!\left[2^{c'_pN\widetilde Q}\mid\text{previous exposure}\right]
 &\le 1+\sum_{j:q_j\le1/2}2^{2c'_pNq_j}2^{-c_pNq_j}\notag\\
 &\le1+2^{-\Omega_{p,\eta}(N/\log N)}.
 \label{eq:path-heavy-one-moment}
\end{align}
The same bound holds at singular starts. There are at most \(2H_0\)
tested columns. Successive conditioning, including between columns of
a pair, yields
\begin{equation}\label{eq:path-heavy-total-moment}
 \E\,2^{c'_pN\sum\widetilde Q}\le1+o(1).
\end{equation}

On \(\mathcal A\), suppose a window node has \(W_R>\eta\).
Iterate Equation~\eqref{eq:path-tail-recursion} backwards from that
node, beginning with \(s=R\) and increasing \(s\) by the width at
each step. Stop once \(s>N^{3/5}\). This takes at most
\(\lceil N^{3/5}\rceil\) steps; even the earliest window node has
\(H_0-5R>\lceil N^{3/5}\rceil\) predecessors for large \(N\).
The final tail vanishes because Lemma~\ref{lem:prime-corank} bounds
every layer height by \(N^{3/5}\).

Number these backwards steps by \(j=1,2,\ldots\). At step \(j\),
the threshold in the recursion is \(R\) plus the preceding widths,
so every charged layer has height at least
\[
 s+1-h\ge R+j-2\ge R-1.
\]
The contribution from one column is therefore at most
\(2Q_{\rm hi}/(R+j-2)\). There are at most two columns per step.
All contributions with \(Q_{\rm hi}<q_{\min}\) total at most
\[
 4q_{\min}\sum_j\frac1{R+j-2}
       \le4q_{\min}(1+\log N)<\eta/2.
\]
The remaining contributions must exceed \(\eta/2\). Their
denominators are at least \(R-1\), and consequently
\begin{equation}\label{eq:path-heavy-forced-sum}
 \sum\widetilde Q\ge\frac{(R-1)\eta}{4},
\end{equation}
where the sum may be extended to all tests on the fixed path.
This same necessary condition holds for a violation at any window
node. Markov's inequality and
Equation~\eqref{eq:path-heavy-total-moment} prove the claim.
\end{proof}

\subsection{The terminal window forces two constraints}

Choose a fixed
\(0<\eta<\min\{1,\zeta/4\}\). The preceding lemma leaves only
windows with \(W_R\le\eta\) at every node. On such a window,
the following deterministic argument is the reason that two full
column costs suffice.

\begin{lemma}[The terminal-window alternative]\label{lem:path-window}
Consider \(5R\) consecutive nonsingular principal extensions of
widths one or two, satisfying their gates. Suppose all their nodes
have \(W_R\le\eta\) and the terminal node has \(W_2>\zeta\).
For all sufficiently large \(N\), either at least two distinct steps
are costly, or a width-two step starts at a node with
\(W_2\le\zeta\).
\end{lemma}

\begin{proof}
At a noncostly step of width \(h\), the recursion gives
\begin{equation}\label{eq:path-noncostly}
 W_s(B)\le W_{s+h}(C)+4\sigma
       \qquad(s\ge h).
\end{equation}
Iterate this inequality from \(s=2\) backwards over any block of
\(R\) noncostly steps. The threshold increases by at least \(R\),
so the endpoint satisfies
\[
 W_2\le\delta,\qquad \delta:=\eta+4R\sigma.
\]
Here and below \(\delta<\zeta\) for large \(N\).

Label the window nodes \(X_0,\ldots,X_{5R}\). If there are no
costly steps, the last \(R\) steps contradict the terminal value.
Suppose there is exactly one costly step, at position \(t\).
There must be fewer than \(R\) steps after it, or the same argument
again applies. Thus \(t\ge4R+1\). Both \(X_{t-2}\) and
\(X_{t-1}\) end blocks of \(R\) noncostly steps inside the window,
and hence both have \(W_2\le\delta\).

Assume there is no width-two step whose starting \(W_2\) is at
most \(\zeta\). The two steps at positions \(t-1,t\) must then
both have width one. Applying Equation~\eqref{eq:path-noncostly}
with \(s=1\) to the preceding step gives
\[
 W_1(X_{t-1})\le W_2(X_{t-2})+4\sigma\le\delta+4\sigma.
\]
Even though step \(t\) is costly, its width is one, so the height
comparison in Lemma~\ref{lem:path-algebra} gives
\[
 W_2(X_t)\le W_1(X_{t-1})\le\delta+4\sigma.
\]
The remaining steps are noncostly. For either width,
Equation~\eqref{eq:path-noncostly} with \(s=2\) and monotonicity
of \(W_s\) in \(s\) show that each increases \(W_2\) by at most
\(4\sigma\). There are fewer than \(R\) such steps, whence
\[
 W_2(X_{5R})\le\eta+8R\sigma<\zeta,
\]
a contradiction. Figure~\ref{fig:terminal-window} records the two
adjacent width-one steps used in this argument.
\end{proof}

\begin{figure}[tbp]
\centering
\begin{tikzpicture}[x=1cm,y=1cm,>=Stealth,font=\small]
 \draw[gray!65] (0,1.75)--(12.8,1.75);
 \draw (0,1.64)--(0,1.86) node[above=4pt] {$X_0$};
 \draw (12.8,1.64)--(12.8,1.86) node[above=4pt] {$X_{5R}$};
 \draw (11,1.64)--(11,1.86) node[above=4pt] {$X_t$};
 \node[above] at (6.4,1.9) {terminal window: $5R$ steps};
 \draw[<->,gray!80!black] (11,1.4)--(12.8,1.4)
       node[midway,below] {$<R$ steps};
 \node[draw,rounded corners=2pt,align=center,inner sep=6pt,
       fill=blue!4] (left) at (2.0,0)
       {$X_{t-2}$\\$W_2\le\delta$};
 \node[draw,rounded corners=2pt,align=center,inner sep=6pt,
       fill=blue!4] (middle) at (5.8,0)
       {$X_{t-1}$\\$W_1\le\delta+4\sigma$};
 \node[draw,rounded corners=2pt,align=center,inner sep=6pt,
       fill=orange!7] (right) at (9.6,0)
       {$X_t$\\$W_2\le\delta+4\sigma$};
 \draw[->,thick] (left)--(middle)
       node[midway,above=15pt,align=center] {noncostly\\width one};
 \draw[->,thick] (middle)--(right)
       node[midway,above=15pt,align=center] {only costly step\\width one};
 \draw[densely dashed,gray!65] (11,1.65)--(right.north);
 \node[below=8pt,align=center] at (left.south)
       {preceding $R$ steps\\are noncostly};
 \node[below=8pt,align=center] at (middle.south)
       {also $W_2\le\delta$ from\\its preceding $R$ steps};
 \draw[->,thick] (right.east)--(12.65,0)
       node[pos=.58,below=5pt,align=center] {noncostly\\steps};
 \node[align=center,anchor=north] at (11.4,-1.0)
       {$W_2(X_{5R})\le\eta+8R\sigma$\\$<\zeta$};
\end{tikzpicture}

\caption{The contradiction setup in Lemma~\ref{lem:path-window}:
suppose the window has only one costly step and no width-two step
starts with \(W_2\le\zeta\). The costly step is within \(R\)
steps of the end. Each of its two preceding nodes ends an \(R\)-step
noncostly block, and the displayed adjacent steps must have width
one. The lower chain enlarges these adjacent steps; horizontal
spacing is schematic. The labels are upper bounds propagated by
the proof, not a plot of a matrix statistic.
Here \(\delta=\eta+4R\sigma\).}
\label{fig:terminal-window}
\end{figure}

\begin{proof}[Proof of Proposition~\ref{prop:column-bound}]
Outside a set of probability \(2^{-\omega(N)}\), every nonsingular
terminal matrix belongs to a required fixed-path event \(\mathcal A\).
By Lemma~\ref{lem:path-high-tail} and Equation~\eqref{eq:path-count},
we may also exclude, at the same scale, all such path events having
\(W_R>\eta\) anywhere in their windows.

For a fixed remaining path with terminal \(W_2>\zeta\),
Lemma~\ref{lem:path-window} forces either two costly positions or
a width-two gate with starting \(W_2\le\zeta\). To make the
conditional probability calculation explicit, let \(\mathcal F_t\)
contain all entries exposed through step \(t\). Define \(D_t\)
to be the event that the starting matrix is in its typical set and
step \(t\) is costly. Its starting-matrix indicator is
\(\mathcal F_{t-1}\)-measurable, and
Lemma~\ref{lem:path-retention-cost} gives
\[
 \E[\ind{D_t}\mid\mathcal F_{t-1}]\le p_N.
\]
For any two positions \(r<t\), including consecutive ones,
\begin{equation}\label{eq:path-two-positions}
 \Prob(D_r\cap D_t)
 =\E\!\left[\ind{D_r}
       \E[\ind{D_t}\mid\mathcal F_{t-1}]\right]
 \le p_N\Prob(D_r)\le p_N^2.
\end{equation}
For a specified width-two position, the event of a typical start
with \(W_2\le\zeta\) and a satisfied gate likewise has probability
at most \(p_N^2\). These bounds require no conditioning on how a
backwards path was selected or on any future path requirement: we
upper-bound by the stated events after dropping all other conditions.

There are only \(O(R^2)\) pairs of positions or individual gate
positions in the window. Taking the union over these and all fixed
paths gives
\[
 \Prob\{B\text{ nonsingular},\ W_2(B)>\zeta\}
 \le2^{-\omega(N)}+2^{o(N)}p_N^2
 =2^{-\omega(N)}+2^{-2N(a-2e_0)+o(N)}.
\]
If \(W_2(B)\le\zeta\), Proposition~\ref{prop:static-column}
gives \(|\mathcal S_{L_0}(B)|\le2^{e_0N}\), apart from another
typical exception. For all sufficiently large \(N\), the displayed
bound and these exceptions are at most \(2^{-N(2a-10e_0)}\),
as required.
\end{proof}
